\documentclass[11pt]{article}
\usepackage[margin=0.85in]{geometry}
\usepackage{amsmath,amssymb,amsthm}
\usepackage[hidelinks]{hyperref}
\newtheorem{theorem}{Theorem}
\title{Two genuine rotors for one preserving basis change\\Conjecture 00000007629}
\author{earthking11}
\date{October 6, 2026}
\begin{document}
\maketitle
\section*{Statement and precise scope}
The conjecture states that rotor realizations of orthogonal basis changes
are unique, with orientation reversal as an exception. We disprove literal
uniqueness using a change that preserves orientation. No identification of
opposite rotors appears in the source; a uniqueness statement in the quotient
by the central sign is a different assertion and is not refuted here.

\begin{theorem}
For the Euclidean plane $V=\mathbb R^2$, $Q(v)=v_0^2+v_1^2$, the actual
spin group in its Clifford algebra contains two distinct elements $1$ and
$-1$, both realizing the identity orthonormal-basis change. The change has
determinant $1$, so orientation reversal does not account for this failure
of uniqueness.
\end{theorem}
\begin{proof}
Let $A=\mathrm{Cl}(V,Q)$ and write $\iota:V\to A$ for the vector inclusion.
The Clifford relation is $\iota(v)^2=Q(v)1$. The standard basis vectors
$e_0,e_1$ have $Q(e_i)=1$ and zero polar pairing when their indices differ.
Thus this is an actual orthonormal basis for the specified form.

Both $\iota(e_0)$ and $\iota(-e_0)$ are units, each its own inverse. They
belong to the vector-unit generators of the Lipschitz group, and
\[
  \iota(e_0)\iota(-e_0)=-\iota(e_0)^2=-1.
\]
Thus $-1$ belongs to the Lipschitz group and is even. Its Clifford
conjugate is itself, so its product with that conjugate is $1$. This proves
pin unitarity, and hence membership in the actual spin group. The unit $1$
also belongs to this group. The real scalar map into $A$ is injective;
since $2\neq0$ in $\mathbb R$, these two elements are distinct.

A rotor $r$ realizes a map $T$ here precisely when
$r\iota(v)\operatorname{star}(r)=\iota(Tv)$ for every vector $v$.
Direct calculation gives
\[
  1\iota(v)\operatorname{star}(1)=\iota(v),\qquad
  (-1)\iota(v)\operatorname{star}(-1)=\iota(v).
\]
Both rotors therefore realize $T=\mathrm{id}_V$. This map fixes every
standard basis vector, preserves $Q$, and has determinant $1$. Since the
two realizations are distinct, there cannot be a unique realizing rotor.
\end{proof}

\section*{Formal verification}
The Lean project uses Mathlib's actual \texttt{CliffordAlgebra} and
\texttt{spinGroup}, not a substitute group. It proves the membership of
$-1$ through generating vector units, pin unitarity and the even subalgebra.
It verifies the squared-norm formula, basis orthonormality, all-vector
conjugation identities, and determinant. The capstone includes these
facts and $\neg\exists!r\in\mathrm{Spin}(Q)$ realizing the identity.

The independently checked project builds with Lean 4.33.1 and Mathlib at
the fixed commit recorded in the manifest. Its printed axiom audits use
only propositional extensionality, classical choice and quotient soundness;
there are no unfinished proofs, native evaluation or additional axioms.
No numerical search or long brute-force computation is used.
\end{document}
